% legends.tex (v10): figure legends (each under 350 words). The years used are stated for every panel.
\section*{Figure legends}

\textbf{Figure 1. ALT observations, inputs and the years they cover.} (a) ESA CCI permafrost probability (mean of 1997--2021) with the three study regions and the number of 1 km cells with ALT labels in 2010--2024. (b) Number of cell-year labels per observation year, 2010--2024, for each region. (c) Periods of the data: labels of each region, the ERA5-Land monthly data used for the year-specific climate input (2010--2024), the static covariates (no year), and the six published products with their available period (light) and the years used for comparison at the label cells (dark); the shaded band marks the analysis period 2010--2024. (d) Definition of the climate input at one Alaska CALM cell (71.3$^\circ$N, 156.6$^\circ$W) in 2014: cumulative thawing degree-days (TDD) from ERA5-Land, the value accumulated to the observation day (year-specific input used for training and scoring; 1 September, the convention for undated CALM end-of-season values), the full-season value of the same year (used for maps) and the 2011--2024 mean (mean climate input).

\textbf{Figure 2. Physics-residual workflow and validation designs.} (a) Structure of the physics-residual model, with the 2014 inputs of Alaska as an example: the year-specific climate variables and the static terrain and soil covariates enter a gradient-boosting model of the residual of the recalibrated Stefan model ($\mathrm{ALT}=E\sqrt{\mathrm{TDD}}$); the observed ALT is used only in training (dashed); the output is the Stefan estimate plus $\lambda=0.5$ times the residual, with the area of applicability (computed with the 2011--2024 mean climate). (b) The four model structures compared: direct ML, physics-input ML (Stefan and Kudryavtsev predictions as features), physics pseudo-labels (training set augmented with Stefan predictions) and physics plus residual ML. (c) The five validation designs shown on the Lena Delta label cells (labels of 2010--2022): random, spatial block (0.5$^\circ$), year hold-out, block-by-year and region hold-out (trained on Alaska and Canada); one example fold is shown for each design, and a cell is drawn as test (black) if any of its cell-years is in the test fold.

\textbf{Figure 3. Year-specific climate input, validation design and the conditions of the residual gain.} All panels use the cell-year labels of 2010--2024 and out-of-fold predictions. (a) Spatial-block RMSE with the mean climate input (2011--2024 mean; open markers) and the year-specific climate input (observation year, accumulated to the observation day; filled markers) for the recalibrated Stefan model, the physics-residual model and direct ML; in this panel only, open and filled markers denote the climate input. (b) RMSE under the five validation designs for the Stefan model with the source coefficient, the recalibrated Stefan model, the physics-residual model and direct ML; year hold-out and block-by-year use the years with at least 20 cell-years; under region hold-out the two Stefan models coincide; values beyond the axis are printed. (c) Difference in spatial-block RMSE between the physics-residual model and the recalibrated Stefan model for six covariate sets. (d) The same difference within regional terciles of eight covariates; cells outlined in black have a 95\% interval excluding zero, hatched cells have fewer than 30 cell-years or fewer than five blocks, and 0.0 marks a difference smaller than 0.05 cm. In (a), (c) and (d) bars and intervals are 95\% block-bootstrap intervals (2,000 resamples of 0.5$^\circ$ blocks); negative differences favour the physics-residual model.

\textbf{Figure 4. Accuracy with few labels and in new regions.} (a) Within-region label ladder: difference in spatial-block RMSE between each method and a Stefan model recalibrated on the same $n$ training labels drawn from the 2010--2024 cell-years (three draws); the band is the 95\% block-bootstrap interval of the physics-residual model. (b) Transfer to a new region: RMSE in the spatial-block folds of each region for models trained on the two other regions plus $n$ labels from the target region. (c) External CALM regions (observations of 2010--2024 outside the study regions): RMSE against the number of labels from the same external region added to the training data, on the \nftrows{} cell-years of seven regions with values of Wei et al., Liu et al. and ESA CCI v5 (2010--2020); the thick light line is the Stefan model recalibrated on the same labels, which nearly coincides with the physics-residual model; horizontal lines show the products, of which Wei et al. and Liu et al. were trained on these sites and ESA CCI v5 was not trained on ALT observations.

\textbf{Figure 5. Comparison with published products in the same cells and years.} (a) RMSE of the three methods (spatial-block out-of-fold predictions) and five products on the cell-years at which Wei et al., Liu et al. and ESA CCI v4 provide a value of the observation year (\nyrscc); filled markers use all common cell-years and open markers only cells more than 25 km from the training sites of Wei et al.; ESA CCI v5 and Aalto et al. are scored on the common cell-years where they also have a value; values beyond the axis are printed. (b) RMSE per observation year on the same cell-years of each year for the physics-residual model, the recalibrated Stefan model, Wei et al. and Liu et al. (to 2020), for years with at least 30 cell-years: Alaska \nyrsywAk, Lena Delta \nyrsywLe{} and Canada \nyrsywCa; grey bars give the number of cell-years per year (log scale).

\textbf{Figure 6. Year-labelled 1 km maps.} (a) Mean ALT of 2011--2024 from the physics-residual model in Alaska, computed from the yearly maps; each yearly map uses the full-season ERA5-Land climate of that year. Dots mark the label cells; lighter colours mark cells outside the area of applicability; grey land lies outside the map domain (ESA CCI permafrost probability below 50\%, water or missing inputs). (b) Anomaly of each year from 2011 to 2024 relative to the 2011--2024 mean, on the same colour range ($\pm$15 cm) for all years. (c) Regional mean anomaly of the predicted maps (lines) for Alaska and the Lena Delta, 2011--2024, and the mean anomaly of the observed cells in years with at least 50 cell-years (circles; Alaska \nanobsyrsAk, Lena Delta \nanobsyrsLe). The 1 km values are interpolated to 250 m for display.
